Quantum-assisted UAV path planning has recently been studied using QAOA-style obstacle-avoidance formulations, but a low-energy binary sample is not necessarily a valid UAV trajectory. This paper presents a feasibility-aware evaluation protocol on three fixed $8\times8$, $L=20$ UAV grid-planning instances with obstacle avoidance, a line-of-sight visibility cost, and a higher-order temporal buffer-risk term. The instances share the start, target, horizon, and objective weights while varying obstacle count and placement; $L$ denotes discrete planning layers, not physical seconds. The model uses binary variables $x_{t,v}$ indicating whether the UAV occupies cell $v$ at time layer $t$ and forms a HUBO objective with one-hot, start, motion, obstacle, visibility, terminal-goal, and sustained near-obstacle proximity terms. The primary six-obstacle instance has 1,280 original binary variables and 118,344 polynomial terms, including 4,392 cubic monomials. We compare A*, RRT-best, a QUAV-style QAOA candidate-path selector, native trajectory-level simulated annealing (SA), the classical \texttt{neal} BQM sampler after HUBO-to-QUBO reduction, and a CP-SAT baseline that exactly solves an integer-scaled hard-feasible objective. CP-SAT returns OPTIMAL on all three instances. On the primary instance, it certifies that the best observed SA path is globally optimal for the implemented scaled objective, with implementation HUBO energy $-20752.43$, soft cost 247.57, path length 14, six turns, and 85\% visibility. Across the primary, two-obstacle, and three-obstacle instances, SA full-task success is 16\%, 84\%, and 66\%, whereas \texttt{neal} full-task success is 0\% in all three. These fixed-instance sensitivity results show why feasibility-aware decoding, exact baselines, and quadratization effects matter; they do not establish quantum speedup or distribution-level performance.
